\PBintroduction{Cinco de Fevereiro}
Outro dia, em meio a essas rotinas modernas, eu estava em uma reunião por videoconferência — daquelas em que precisamos manter a câmera aberta o tempo todo.

Nesse dia, porém, a ferramenta apresentou um erro. Precisei sair e acessar novamente. 

Já quase no final da reunião, por conta desse vai e volta, o fundo que costumo utilizar acabou se desconfigurando.

Foi então que me deparei com uma parede branca atrás de mim.

Demorei a perceber.

E isso chama a atenção, porque meu fundo padrão era uma imagem preta e amarela, com frases, nomes e acrônimos da área de TI.

Mas, por alguns minutos, eu simplesmente não vi.

Isso aconteceu no dia cinco de fevereiro de 2026. 

Uma quinta-feira.

Ocorre que há algum tempo — cerca de sete meses antes — vinha fazendo terapia, sempre às quartas-feiras. 

E, pelo menos nos últimos sete encontros, uma constatação tem se repetido: eu não tenho um propósito de vida.

Não há nada que me faça pensar: “vou viver para isso…”, “quero fazer aquilo…”, ou ainda “meu sonho é…”.

A partir disso, diversos pensamentos começaram a surgir.

Diversos “porquês”.

Por que uma parede branca? 

Porque não tem quadros? 

Porque não tem cor?

Reflexões que, até então, eu nunca tinha feito. 

Tais como, não tenho conquista que vale a pena compartilhar? 

Não tenho nada do que me orgulhar? 

Preciso ter algo na parede?

Talvez precise levar vocês para uma jornada.

E que fique claro. 

Isso não é autoajuda. 

Não é autopiedade. Não é conselho.

É apenas um relato.

Uma história de vida — como a de tantas outras pessoas.

Uma trajetória de quase meio século.

E se por acaso se conectar com a sua realidade, se trata de mera coincidência. 

Sei que as vidas são únicas, trajetórias são únicas. 

%\PBreflection{A história não se repete, mas que ela rima, ela rima}